% Irodov (user-provided Arihant PDF), printed pp.29--30, PDF pp.38--39.
% Original problem numbers are preserved. Calculations independently expanded.
\exproblem{171}{29}
一炮弹以 $v=500\,\mathrm{m/s}$ 飞行，突然炸成质量相等的三块。
爆炸后三块的总动能是爆炸前的 $\eta=1.5$ 倍。
忽略爆炸瞬间外力冲量，问其中一块的最大可能速率。
\exsolution
设原总质量为 $M$，其中一块速度为 $\bm u$，
另两块速度为 $\bm v_2,\bm v_3$，爆炸前速度为 $\bm v$。
动量守恒给出
$\bm u+\bm v_2+\bm v_3=3\bm v$。
对于固定的 $\bm u$，另外两块的动能和在
$\bm v_2=\bm v_3=(3\bm v-\bm u)/2$ 时最小，
因为平方和不小于两倍平均矢量模平方。于是
\[
 K_f\ge\frac M6\left[
 u^2+\frac12|3\bm v-\bm u|^2\right]
 =\frac M2v^2+\frac M4|\bm u-\bm v|^2 .
\]
结合 $K_f=\eta Mv^2/2$ 得
$|\bm u-\bm v|^2\le2(\eta-1)v^2$。
三角不等式及其可取等条件给出
\[
\boxed{u_{\max}=v\bigl[1+\sqrt{2(\eta-1)}\bigr]
=1000\,\mathrm{m/s}}.
\]
最大值在该碎片与原速度同向、另两碎片同速时取得。
\excheck 对 $\eta=1$，没有额外爆炸动能时上界回到 $v$；
速度是矢量，单用一维能量守恒而不考虑动量约束会高估结果。

\exproblem{172}{29}
质量相等的两质点发生一维正碰。质点1碰前速率
$v=10\,\mathrm{m/s}$，质点2静止；碰后总动能损失
初始动能的 $\eta=1.0\%$。求碰后质点1的速度（含方向）。
\exsolution
取入射方向为正，碰后速度为 $u_1,u_2$。
动量守恒给出 $u_1+u_2=v$；碰后总动能为
$(1-\eta)mv^2/2$，因此
\[
u_1^2+(v-u_1)^2=(1-\eta)v^2,
\]
\[
u_1=\frac v2\left[1\pm\sqrt{1-2\eta}\right].
\]
碰后两球必须分离，因而 $u_2\ge u_1$；
选用减号分支，得到
\[
\boxed{u_1=\frac v2(1-\sqrt{1-2\eta})
\simeq0.0503\,\mathrm{m/s}},
\]
方向仍与原入射方向相同；$u_2\simeq9.9497\,\mathrm{m/s}$。
\excheck 加号分支意味着后一球比前一球慢，碰后仍会相互接近，
不代表一次完成的分离碰撞；题设要求 $0\le\eta\le1/2$。

\exproblem{173}{29}
质量 $m$ 的入射粒子碰前速率为 $v$，与静止质量
$M$ 的粒子碰撞后偏转 $90^\circ$；靶粒子的反冲方向与原
入射方向成 $\theta=30^\circ$，且 $M/m=5.0$。
求系统总动能的百分变化。
\exsolution
取原入射速度沿 $x$ 轴，入射粒子末速度沿 $+y$，
靶粒子反冲必在 $-y$ 一侧。设末速分别为 $u,V$，
两个方向的动量守恒为
\[
mv=MV\cos\theta,\qquad mu=MV\sin\theta.
\]
因此 $V=mv/(M\cos\theta)$、
$u=v\tan\theta$。末初动能比
\[
\frac{K_f}{K_i}
=\frac{mu^2+MV^2}{mv^2}
=\tan^2\theta+\frac mM\sec^2\theta.
\]
所以
\[
\boxed{\frac{\Delta K}{K_i}
=\left(1+\frac mM\right)\tan^2\theta+
\frac mM-1=-0.40}.
\]
\excheck 系统总动能减少 $40\%$，不是增加；此推导只要求
碰撞期间总动量守恒，不需要假定弹性碰撞。

\exproblem{174}{29}
孤立系统内两个质点质量为 $m_1,m_2$，
原参考系中的速度 $\bm v_1,\bm v_2$ 互相垂直。
求两质点在质心参考系中的动量及该系总动能。
\exsolution
质心速度
$\bm V=(m_1\bm v_1+m_2\bm v_2)/(m_1+m_2)$。
定义约化质量 $\mu=m_1m_2/(m_1+m_2)$，则
\[
\widetilde{\bm p}_1
=m_1(\bm v_1-\bm V)=\mu(\bm v_1-\bm v_2),
\quad
\widetilde{\bm p}_2=-\widetilde{\bm p}_1.
\]
由 $\bm v_1\cdot\bm v_2=0$，
\[
\boxed{|\widetilde{\bm p}_1|
=|\widetilde{\bm p}_2|=\mu\sqrt{v_1^2+v_2^2}},
\qquad
\boxed{\widetilde K
=\frac{\mu}{2}|\bm v_1-\bm v_2|^2
=\frac{\mu}{2}(v_1^2+v_2^2)}.
\]
\excheck 质心系总动量为零，但总动能通常不为零；
这里没有把两个粒子的\emph{原}动量当作质心系动量。

\exproblem{175}{29}
质量 $m_1>m_2$ 的粒子撞击静止质量 $m_2$ 的粒子，
碰撞完全弹性。求入射粒子在碰后相对原入射方向的
最大偏转角。
\exsolution
设入射速率 $v$，碰后入射粒子速度矢量 $\bm u$。
消去靶粒子的末速度
$\bm w=(m_1/m_2)(\bm v-\bm u)$，
由动能守恒得
\[
(m_1+m_2)u^2-2m_1vu\cos\theta+(m_1-m_2)v^2=0.
\]
把它视为关于 $u$ 的二次方程，要求存在实正根，
判别式非负：
\[
m_1^2\cos^2\theta-(m_1+m_2)(m_1-m_2)\ge0,
\quad \sin^2\theta\le\frac{m_2^2}{m_1^2}.
\]
所以
\[
\boxed{\theta_{\max}=\arcsin\frac{m_2}{m_1}}.
\]
\excheck $m_1\gg m_2$ 时偏转角很小；
当 $m_1\downarrow m_2$ 时最大偏转趋于 $90^\circ$。

\exproblem{176}{29}
三块质量相同的光滑圆盘 $A,B,C$ 在光滑水平面上。
$B,C$ 静止、两盘心间距等于盘直径的 $\eta$ 倍；
$A$ 沿两盘心连线的垂直平分线前进，速率 $v$，
同时与 $B,C$ 完全弹性碰撞（原图1.45）。
求 $A$ 碰后速度与反弹、停止、继续前进的判据。
\exsolution
盘半径设为 $R$，同时接触要求 $1\le\eta<2$。
碰撞时 $\overline{AB},\overline{AC}$ 与入射方向
的夹角为 $\phi$，几何给出
$\sin\phi=(2R\eta/2)/(2R)=\eta/2$，所以
$\cos^2\phi=1-\eta^2/4$。
设对 $B,C$ 的法向冲量大小均为 $J$，对称性使
$A$ 碰后仍仅沿原直线，速度有符号值 $u$：
\[
m(v-u)=2J\cos\phi.
\]
两目标盘碰后速度大小都为 $J/m$；
由于碰撞弹性，
\[
mv^2=mu^2+2J^2/m.
\]
解以上两式、舍去无碰撞的零冲量分支，得到
\[
\boxed{u=v\frac{\eta^2-2}{6-\eta^2}}.
\]
因此 $1\le\eta<\sqrt2$ 时 $A$ 反弹；
$\eta=\sqrt2$ 时停止；
$\sqrt2<\eta<2$ 时继续向前。
\excheck 这里的 $u$ 是\emph{有符号速度}；图形的双接触
只在指定中心距区间有意义，不能延伸到 $\eta>2$。

\exproblem{177}{29--30}
一个分子与另一个同质量、碰前静止的分子碰撞。
证明：（a）碰撞理想弹性时，两粒子的末速度方向夹角
为 $90^\circ$；（b）若动能损失为正，夹角不再是 $90^\circ$。
\exsolution
两粒子质量均为 $m$，碰前入射速度为 $\bm v$，
末速度为 $\bm u,\bm w$。动量守恒
$\bm v=\bm u+\bm w$，故
\[
v^2=u^2+w^2+2\bm u\cdot\bm w.
\]
若完全弹性，动能守恒
$v^2=u^2+w^2$，所以
$2\bm u\cdot\bm w=0$；当二者末速都非零时，
\[
\boxed{\angle(\bm u,\bm w)=90^\circ}.
\]
若发生非弹性碰撞且动能损失 $\Delta E>0$，
则 $v^2>u^2+w^2$，故
\[
2\bm u\cdot\bm w
=v^2-u^2-w^2=2\Delta E/m>0,
\]
夹角为锐角。
\excheck 如有一粒子恰停下，其速度方向未定义；
“夹角为 $90^\circ$”指两末速度非零的通常散射情形。

\exproblem{178}{30}
火箭在某时刻质量为 $m$、加速度为 $\bm a$，
外力合力为 $\bm F$。火箭以恒定质量流率
$\mu>0$（$\mathrm{kg/s}$）喷出气体，
气体相对火箭的速度矢量为 $\bm u$。
证明火箭方程 $m\bm a=\bm F-\mu\bm u$。
\exsolution
取时间间隔 $\mathrm dt$，火箭初始速度 $\bm v$，
其质量减少 $\mu\,\mathrm dt$；
喷出的气体质量为 $\mu\,\mathrm dt$，
在惯性系中速度为 $\bm v+\bm u$。
动量变化保留到 $\mathrm dt$ 一阶：
\[
\begin{aligned}
\Delta\bm P
&=(m-\mu\mathrm dt)(\bm v+\mathrm d\bm v)
 +\mu\mathrm dt(\bm v+\bm u)-m\bm v\\
&=m\,\mathrm d\bm v+\mu\bm u\,\mathrm dt
+O(\mathrm dt^2).
\end{aligned}
\]
由外力冲量 $\bm F\,\mathrm dt=\Delta\bm P$，
\[
\boxed{m\bm a=\bm F-\mu\bm u}.
\]
\excheck 此处 $\bm u$ 是\emph{喷气相对火箭的矢量速度}，
通常指向火箭后方；反冲推力方向为 $-\mu\bm u$。

\exproblem{179}{30}
火箭初始静止、质量 $m_0$，在无外力情况下
以相对火箭不变的喷气速度 $\bm u$ 排气。
求火箭质量变为 $m$ 时的速度。
\exsolution
令 $\mu=-\mathrm dm/\mathrm dt>0$，
用上一题方程且 $\bm F=\bm0$：
\[
m\frac{\mathrm d\bm v}{\mathrm dt}
=-\mu\bm u=\bm u\,\frac{\mathrm dm}{\mathrm dt}.
\]
两端分离变量并由初态积分：
\[
\int_{\bm0}^{\bm v}\mathrm d\bm v
=\bm u\int_{m_0}^{m}\frac{\mathrm dm'}{m'},
\]
所以
\[
\boxed{\bm v=\bm u\ln\frac m{m_0}
=-\bm u\ln\frac{m_0}{m}}.
\]
\excheck $m<m_0$，若 $\bm u$ 指向火箭后方，
速度 $\bm v$ 指向前方；其大小是
$|\bm u|\ln(m_0/m)$，与喷气质量流率无关。

\exproblem{180}{30}
在无外力作用下，火箭保持大小为 $a>0$ 的恒加速度，
喷气相对火箭的反向速率固定为 $u>0$，
初质量为 $m_0$。求质量 $m(t)$。
\exsolution
设推进方向为正，气体相对火箭速度为 $-u$。
火箭方程 $ma=-\mu(-u)=\mu u$，
其中 $\mu=-\dot m$，故
\[
ma=-u\frac{\mathrm dm}{\mathrm dt},
\qquad\frac{\mathrm dm}{m}=-\frac au\,\mathrm dt.
\]
由 $m(0)=m_0$ 得
\[
\boxed{m(t)=m_0e^{-at/u}}.
\]
\excheck 指数 $at/u$ 无量纲；质量单调减少且始终为正。
恒加速度要求质量流率按 $\mu(t)=am(t)/u$
随时间衰减，并非恒定流率。
